\chapter[Nierówności wariacyjne i przeszkoda]{Nierówności wariacyjne, zagadnienie przeszkody i swobodny brzeg}
\label{chap:pde29}

\section{Motywacja: energia z ograniczeniem jednostronnym}
W poprzednich rozdziałach minimalizowaliśmy funkcjonały energii na całych
przestrzeniach Sobolewa. W wielu zagadnieniach mechanicznych funkcja
$u$ nie może jednak przekroczyć przeszkody w ustalonym kierunku.
Na przykład odkształcenie membrany nie może być mniejsze niż wysokość
sztywnej podpory $\psi$. W takim przypadku warunek Eulera--Lagrange'a
nie musi być równaniem w całym obszarze. Na zbiorze kontaktu pojawia
się dodatkowa reakcja, a sam zbiór kontaktu nie jest znany z góry.

Niech $\Omega\subset\R^d$ będzie ograniczonym obszarem Lipschitzowskim.
Pracujemy na rzeczywistych przestrzeniach funkcyjnych.
Oznaczamy $V=H_0^1(\Omega)$, z iloczynem skalarnym
$(v,w)_V=\int_\Omega\nabla v\cdot\nabla w\,\mathrm d\lambda_d(x)$
i normą $\|v\|_V=\|\nabla v\|_{L^2(\Omega)}$.
Nierówność Poincarégo~\ref{cor:pde10-Poincare}
daje równoważność tej normy z normą $H^1$, więc
$V$ jest przestrzenią Hilberta.
Przez $V'=H^{-1}(\Omega)$ oznaczamy jej przestrzeń
dualną z normą względem $\|\cdot\|_V$.
Wszystkie relacje $v\ge w$ rozumiemy prawie wszędzie.
Przyjmujemy $\psi\in H^1(\Omega)$ i $\psi^+\in V$,
gdzie $\psi^+=\max(\psi,0)$ należy do $H^1$
na mocy Wniosku~\ref{cor:pde9-truncation}.
Ostatni warunek gwarantuje istnienie funkcji dopuszczalnej.
Definiujemy
\[
 K_\psi=\{v\in V:v\geq\psi\ \text{prawie wszędzie w }\Omega\}.
\]
Zbiór ten jest niepusty, bo $\psi^+\in K_\psi$, oraz wypukły.
Przez $f\in V'$ oznaczamy obciążenie, a przez
$\langle f,v\rangle$ wartość funkcjonału $f$ na $v$.

\begin{lemat}[Domkniętość zbioru funkcji dopuszczalnych]
\label{lem:pde29-K}
Dla ograniczonego obszaru Lipschitzowskiego $\Omega\subset\R^d$
i dowolnej funkcji $\psi\in H^1(\Omega)$ spełniającej
$\psi^+\in H_0^1(\Omega)$ zbiór $K_\psi$ jest niepusty,
wypukły i słabo domknięty w $H_0^1(\Omega)$.
\end{lemat}
\begin{proof}
Niepustość wynika z $\psi^+\in K_\psi$.
Dla $v,w\in K_\psi$ i $\theta\in[0,1]$ mamy
$\theta v+(1-\theta)w\ge\psi$ prawie wszędzie.
Dla $v\in V$ zachodzi równoważność
\[
 v\in K_\psi\quad\Longleftrightarrow\quad
 \int_\Omega(v-\psi)g\,\mathrm d\lambda_d\ge0
 \quad\text{dla każdego }g\in L^2(\Omega),\ g\ge0.
\]
Implikacja w prawo wynika z nieujemności czynników.
W drugą stronę wybieramy $g=(\psi-v)^+$:
całka wynosi $-\|(\psi-v)^+\|_2^2$, więc $v\ge\psi$.
Dla ustalonego $g$ funkcjonał
$v\mapsto\int_\Omega vg\,\mathrm d\lambda_d$
jest ciągły w $V$, na mocy nierówności Poincarégo
i nierówności Cauchy'ego--Schwarza.
Zatem $K_\psi$ jest przecięciem słabo domkniętych
półprzestrzeni, a więc jest słabo domknięty,
nie tylko słabo ciągowo domknięty.
\end{proof}

\section{Rozwiązanie wariacyjne i jego jednoznaczność}
Niech $A(x)=(a_{ij}(x))_{i,j=1}^d$ będzie mierzalną macierzą
symetryczną, $A\in L^\infty(\Omega;\R^{d\times d})$, oraz
$\lambda,\Lambda>0$ takie, że dla prawie każdego $x\in\Omega$
i każdego $\xi\in\R^d$ zachodzi
$\lambda|\xi|^2\le\xi^{\mathsf T}A(x)\xi\le\Lambda|\xi|^2$.
Wprowadzamy formę $a(u,v)=\int_\Omega A\nabla u\cdot\nabla v\,
\mathrm d\lambda_d(x)$ i funkcjonał
$J_f(v)=\frac12a(v,v)-\langle f,v\rangle$.
Z symetrii $A(x)$ i oszacowań jej wartości własnych
wynika $\|A(x)\|_{\mathrm{op}}\le\Lambda$ prawie wszędzie.
Nierówność Cauchy'ego--Schwarza daje
$|a(u,v)|\le\Lambda\|u\|_V\|v\|_V$,
a eliptyczność $a(v,v)\ge\lambda\|v\|_V^2$.
Forma $a$ jest iloczynem skalarnym, którego norma
$\|v\|_a=\sqrt{a(v,v)}$ spełnia
\[
 \sqrt\lambda\|v\|_V\le\|v\|_a\le
 \sqrt\Lambda\|v\|_V.
\]
Jest więc zupełna i wyznacza tę samą topologię
słabą, gdyż obie normy mają te same ciągłe
funkcjonały liniowe.

\begin{twierdzenie}[Istnienie i jednoznaczność w zagadnieniu przeszkody]
\label{thm:pde29-existence}
Niech $\Omega$, $A$, $\psi$ spełniają powyższe założenia,
a $f\in H^{-1}(\Omega)$. Istnieje dokładnie jedna funkcja
$u\in K_\psi$, dla której
\begin{equation}\label{eq:pde29-VI}
 a(u,v-u)\ge\langle f,v-u\rangle
 \qquad\text{dla każdego }v\in K_\psi.
\end{equation}
Jest ona jedynym punktem, w którym $J_f$ osiąga minimum na $K_\psi$.
\end{twierdzenie}
\begin{proof}
Z eliptyczności, definicji normy dualnej
i nierówności Younga mamy
\[
 J_f(v)\ge\tfrac\lambda2\|v\|_V^2-\|f\|_{V'}\|v\|_V
 \ge\tfrac\lambda4\|v\|_V^2-\lambda^{-1}\|f\|_{V'}^2.
\]
Niech $m=\inf_{K_\psi}J_f$. Z niepustości zbioru
i powyższego oszacowania $-\infty<m\le J_f(\psi^+)<\infty$.
Wybieramy $u_n\in K_\psi$ z $J_f(u_n)\le m+1/n$.
Ciąg jest ograniczony w $V$; refleksywność i twierdzenie
Eberleina--Šmuliana pozwalają wydzielić podciąg
słabo zbieżny do $u\in V$.
Z Lematu~\ref{lem:pde29-K} mamy $u\in K_\psi$.
Norma $\|\cdot\|_a$ jest słabo półciągła z dołu,
bo $\|u\|_a=\sup_{\|w\|_a\le1}|a(u,w)|$
i każdy funkcjonał $a(\cdot,w)$ jest słabo ciągły.
Składnik liniowy jest słabo ciągły, więc
$m\le J_f(u)\le\liminf_n J_f(u_n)=m$.
Otrzymujemy minimizator.

Dla $v\in K_\psi$ odcinek $u+t(v-u)$ należy do $K_\psi$
dla $t\in[0,1]$. Minimalność daje
$0\le[J_f(u+t(v-u))-J_f(u)]/t
 =a(u,v-u)-\langle f,v-u\rangle+
 \frac t2a(v-u,v-u)$. Przechodząc z $t$ do zera od strony dodatniej,
otrzymujemy~\eqref{eq:pde29-VI}. Odwrotnie, dla funkcji
spełniającej tę nierówność
$J_f(v)-J_f(u)=a(u,v-u)-\langle f,v-u\rangle
+\frac12a(v-u,v-u)\ge0$.
Dokładniej, nierówność wariacyjna i eliptyczność dają
\[
 J_f(v)-J_f(u)\ge\tfrac12\|v-u\|_a^2
 \ge\tfrac\lambda2\|v-u\|_V^2.
\]
Każdy ciąg minimalizujący zbiega więc silnie do $u$
w normie $V$.
Jeżeli $u_1,u_2$ spełniają nierówność, testujemy pierwszą przez
$v=u_2$, drugą przez $v=u_1$ i dodajemy. Otrzymujemy
$a(u_1-u_2,u_1-u_2)\le0$, skąd $u_1=u_2$.
\end{proof}


\begin{stwierdzenie}[Zagadnienie przeszkody jako rzut energetyczny]
\label{prop:pde-v44-obstacle-projection}
Niech $V=H_0^1(\Omega)$, $a$ będzie symetryczną,
ciągłą i koercywną formą z twierdzenia~\ref{thm:pde29-existence},
a $K_\psi\subset V$ niepustym, domkniętym i wypukłym
zbiorem dopuszczalnym. Niech $z\in V$ będzie jedynym
rozwiązaniem równania $a(z,v)=\langle f,v\rangle$ dla
każdego $v\in V$. Wtedy rozwiązanie przeszkody $u$
jest rzutem metrycznym $z$ na $K_\psi$ w normie
$\|v\|_a=a(v,v)^{1/2}$. Co więcej, dla ustalonego
$K_\psi$ i dwóch danych $f_i\in V'$ zachodzi
\begin{equation}\label{eq:pde-v44-projection-nonexpansive}
 \|u_1-u_2\|_a\le \|z_1-z_2\|_a
 \le\lambda^{-1/2}\|f_1-f_2\|_{V'},
\end{equation}
gdzie $z_i$ rozwiązuje równanie liniowe dla $f_i$,
a $u_i$ odpowiednie zagadnienie przeszkody.
\end{stwierdzenie}
\begin{proof}
Istnienie i jednoznaczność $z$ wynikają z twierdzenia Laxa--Milgrama;
forma $a$ jest iloczynem skalarnym równoważnym
standardowemu na $V$. Równanie dla $z$ daje
\[
 J_f(v)=\tfrac12\|v-z\|_a^2-\tfrac12\|z\|_a^2.
\]
Minimalizacja energii jest zatem dokładnie
minimalizacją odległości od $z$. Z warunku
$a(u,v-u)\ge\langle f,v-u\rangle=a(z,v-u)$
dla $v\in K_\psi$ otrzymujemy
$a(u-z,v-u)\ge0$. Jest to dokładnie
charakteryzacja rzutu na domknięty zbiór wypukły
w przestrzeni Hilberta $(V,a)$.
Dla dwóch rzutów wykorzystujemy w nierównościach
wariacyjnych testy $v=u_2$ i $v=u_1$, a następnie
dodajemy. Dla $w=u_1-u_2$ dostajemy
$a(w,w)\le a(z_1-z_2,w)$; nierówność Cauchy'ego--Schwarza
daje pierwsze oszacowanie, po podzieleniu
przez $\|w\|_a$ dla $w\ne0$.
Dla $\delta z=z_1-z_2$ testujemy równanie
liniowe przez $\delta z$:
\[
 \|\delta z\|_a^2=\langle f_1-f_2,\delta z\rangle
 \le\lambda^{-1/2}\|f_1-f_2\|_{V'}\|\delta z\|_a.
\]
Dla $\delta z\ne0$ dzielimy przez normę, a dla
$\delta z=0$ oszacowanie jest natychmiastowe.
\end{proof}


\begin{stwierdzenie}[Stabilność względem obciążenia]
\label{prop:pde29-stable}
Przy ustalonych $\Omega$, $A$, $\psi$ z
 twierdzenia~\ref{thm:pde29-existence}, jeżeli $u_i$ odpowiadają
$f_i\in V'$ ($i=1,2$), to
$\|u_1-u_2\|_V\le\lambda^{-1}\|f_1-f_2\|_{V'}$.
\end{stwierdzenie}
\begin{proof}
W obu nierównościach wariacyjnych wybieramy jako test funkcję
będącą drugim rozwiązaniem. Po dodaniu i zastosowaniu
eliptyczności otrzymujemy
$\lambda\|u_1-u_2\|_V^2\le
\langle f_1-f_2,u_1-u_2\rangle$.
Dalej stosujemy nierówność dualną; przypadek $u_1=u_2$
nie wymaga dzielenia przez normę.
\end{proof}

\begin{twierdzenie}[Zasada porównania dla przeszkód]
\label{thm:pde29-order}
Niech $A$, $\Omega$ spełniają powyższe założenia,
a $\psi_i\in H^1(\Omega)$, $\psi_i^+\in V$ oraz
$f_i\in L^2(\Omega)$, $i=1,2$. Jeśli
$\psi_1\le\psi_2$ i $f_1\le f_2$ prawie wszędzie,
to odpowiadające rozwiązania spełniają $u_1\le u_2$
prawie wszędzie w $\Omega$.
\end{twierdzenie}
\begin{proof}
Niech $w=(u_1-u_2)^+$.
Z Wniosku~\ref{cor:pde9-truncation} mamy $w\in H^1$
oraz $\nabla w=\chi_{\{u_1>u_2\}}\nabla(u_1-u_2)$.
Aby wykazać $w\in V$, przybliżamy różnicę
$u_1-u_2$ przez funkcje $g_n\in C_c^\infty(\Omega)$
w normie $H^1$. Funkcje $g_n^+$ mają zwarty
nośnik w $\Omega$ i należą do $V$ przez wygładzanie.
Ich normy $H^1$ są ograniczone, a
$\|g_n^+-w\|_2\le\|g_n-(u_1-u_2)\|_2\to0$.
Słaba domkniętość $V$ w $H^1$ identyfikuje
granicę jako element $V$.
Ponieważ
$\min(u_1,u_2)\ge\psi_1$ i
$\max(u_1,u_2)\ge\psi_2$,
w nierówności dla $u_1$ wolno wziąć $v=u_1-w$,
a w nierówności dla $u_2$ funkcję $v=u_2+w$.
Daje to $a(u_1,w)\le\int_\Omega f_1w\,\mathrm d\lambda_d(x)$
oraz $a(u_2,w)\ge\int_\Omega f_2w\,\mathrm d\lambda_d(x)$.
Po odjęciu i wykorzystaniu $\nabla w=\chi_{\{u_1>u_2\}}
\nabla(u_1-u_2)$ otrzymujemy
$\lambda\|w\|_V^2\le\int_\Omega(f_1-f_2)w\,\mathrm d\lambda_d(x)
\le0$. Stąd $w=0$.
\end{proof}


\begin{twierdzenie}[Ciągła zależność od przeszkody w normie energetycznej]
\label{thm:pde29-obstacle-stability}
Niech $\Omega\subset\R^d$ będzie ograniczonym obszarem
Lipschitzowskim, $V=H_0^1(\Omega)$, a $A$ niech będzie
symetryczna, mierzalna i niech spełnia warunki
$\lambda|\xi|^2\le\xi^{\mathsf T}A(x)\xi\le\Lambda|\xi|^2$
dla $\lambda,\Lambda>0$. Niech $f\in V'$ będzie ustalone.
Dla $j=1,2$ wybierzmy przeszkodę $\psi_j\in V$
i oznaczmy przez $u_j\in K_{\psi_j}$ rozwiązanie
nierówności wariacyjnej~\eqref{eq:pde29-VI}.
Wtedy
\begin{equation}\label{eq:pde29-obstacle-stability}
 \|u_1-u_2\|_V\le\sqrt{\frac{\Lambda}{\lambda}}
 \|\psi_1-\psi_2\|_V.
\end{equation}
\end{twierdzenie}
\begin{proof}
Niech $h=\psi_1-\psi_2\in V$.
Ponieważ $u_2\ge\psi_2$, mamy $u_2+h\ge\psi_1$;
podobnie $u_1-h\ge\psi_2$. Są to zatem poprawne
funkcje testowe w obu nierównościach wariacyjnych.
Pierwszą testujemy przez $u_2+h$, a drugą przez $u_1-h$.
Po dodaniu, przy $w=u_1-u_2$, otrzymujemy
\[
 a(w,w-h)\le\langle f,w-h\rangle
                 -\langle f,w-h\rangle=0.
\]
Stąd $a(w,w)\le a(w,h)$.
Nierówność Cauchy'ego--Schwarza w iloczynie
skalarnym $a$ daje $\|w\|_a^2\le\|w\|_a\|h\|_a$,
czyli $\|w\|_a\le\|h\|_a$ (dla $w=0$ nie dzielimy).
Z równoważności norm dostajemy
\[
 \|w\|_V\le\lambda^{-1/2}\|w\|_a
 \le\lambda^{-1/2}\|h\|_a
 \le\sqrt{\Lambda/\lambda}\|h\|_V.
\]
Jest to~\eqref{eq:pde29-obstacle-stability}.
\end{proof}

\begin{uwaga}[Ograniczenie założenia o przeszkodach]
Twierdzenie dotyczy przeszkód $\psi_j\in H_0^1(\Omega)$,
a nie tylko warunku $\psi_j^+\in H_0^1(\Omega)$.
Jest to istotne: przesunięte funkcje testowe
$u_2+\psi_1-\psi_2$ i $u_1-\psi_1+\psi_2$
muszą należeć do $H_0^1(\Omega)$.
\end{uwaga}

\section{Reakcja przeszkody i warunki komplementarności}
W tej sekcji rozważamy $A=I$. Oznaczamy przez $-\Delta u$
dystrybucję określoną wzorem
$\langle-\Delta u,\varphi\rangle
=\int_\Omega\nabla u\cdot\nabla\varphi\,\mathrm d\lambda_d(x)$.

\begin{stwierdzenie}[Dodatnia reakcja jako miara]
\label{prop:pde29-lambda-measure}
Niech $u\in K_\psi$ rozwiązuje~\eqref{eq:pde29-VI}
dla $A=I$ i $f\in V'$. Dystrybucja
$\mu=-\Delta u-f$ jest nieujemna, czyli
$\langle\mu,\varphi\rangle\ge0$ dla każdego
$\varphi\in C_c^\infty(\Omega)$, $\varphi\ge0$.
Wobec tego $\mu$ jest nieujemną miarą Radona w $\Omega$.
\end{stwierdzenie}
\begin{proof}
Jeśli $\varphi\in V$ i $\varphi\ge0$ prawie wszędzie,
to $u+\varphi\in K_\psi$,
więc~\eqref{eq:pde29-VI} daje
$\langle-\Delta u-f,\varphi\rangle\ge0$.
W szczególności $\mu$ jest dodatnią dystrybucją
i ciągłym funkcjonałem na $V$.
Dla dowolnego zwartego $K\Subset\Omega$ wybierzmy
$\eta\in C_c^\infty(\Omega)$, $\eta\ge0$, $\eta=1$
w otoczeniu $K$. Jeżeli $\varphi\in C_c^\infty(\Omega)$ i
$\operatorname{supp}\varphi\subset K$,
to $-\|\varphi\|_\infty\eta\le\varphi\le
\|\varphi\|_\infty\eta$. Dodatniość daje
$|\langle\mu,\varphi\rangle|\le
\|\varphi\|_\infty\langle\mu,\eta\rangle$.
Dla $g\in C_c(\Omega)$ wybieramy gładkie
przybliżenia $g_n$ z nośnikami w jednym zwartym
podzbiorze $\Omega$, zbiegające jednostajnie do $g$.
Można je otrzymać przez przedłużenie przez zero
i splot z jądrem wygładzającym o małym nośniku.
Powyższe oszacowanie zapewnia istnienie granicy
$L(g)=\lim_n\langle\mu,g_n\rangle$ oraz jej
niezależność od przybliżeń (dla dwóch ciągów
stosujemy oszacowanie na sumie ich nośników).
Funkcjonał $L$ jest liniowy i dodatni:
dla $g\ge0$ wybieramy nieujemne jądro splotowe.
Twierdzenie Riesza o reprezentacji dodatniego
funkcjonału na $C_c(\Omega)$ daje jedyną
nieujemną miarę Radona reprezentującą $L$.
Jej działanie na funkcjach testowych jest
równe działaniu pierwotnej dystrybucji.
Nie wnioskujemy, że miara ma skończoną masę
na całej $\Omega$. Nie utożsamiamy też parowania
w $V',V$ z całką dowolnego reprezentanta
funkcji Sobolewa względem $\mu$; reprezentanci
równi prawie wszędzie względem $\lambda_d$
mogą różnić się na zbiorze o dodatniej mierze $\mu$.
\end{proof}

\begin{uwaga}[Komplementarność w parowaniu dualnym]
Jeżeli dodatkowo $\psi\in V$, to można testować
nierówność wariacyjną przez $v=\psi$.
Otrzymujemy $\langle\mu,u-\psi\rangle\le0$.
Ponieważ $u-\psi\ge0$ i należy do $V$, dodatniość
funkcjonału daje nierówność przeciwną, a zatem
$\langle\mu,u-\psi\rangle=0$.
Odwrotnie, jeśli $u\in K_\psi$, $\mu=-\Delta u-f$
jest dodatnia na stożku nieujemnych elementów $V$
i zachodzi ta równość, to dla każdego $v\in K_\psi$
mamy $\langle\mu,v-u\rangle
=\langle\mu,v-\psi\rangle\ge0$, czyli~\eqref{eq:pde29-VI}.
Przy przeszkodzie spoza $V$ takie parowanie
nie jest na ogół określone.
\end{uwaga}

\begin{lemat}[Pochodne funkcji Sobolewa na zbiorze jej zer]
\label{lem:pde29-zero-level}
Niech $\Omega\subset\R^d$ będzie zbiorem otwartym.
Jeżeli $z\in H^1_{\mathrm{loc}}(\Omega)$ i
$E=\{x\in\Omega:z(x)=0\}$, to $\nabla z=0$
prawie wszędzie na $E$. Jeśli ponadto $z\in H^2_{\mathrm{loc}}(\Omega)$,
to $D^2z=0$ prawie wszędzie na $E$.
\end{lemat}
\begin{proof}
Pierwszą część można otrzymać z Lematu
~\ref{lem:pde9-levelsets}, stosując go do ograniczeń
$z$ na relatywnie zwartych kostkach zawartych w $\Omega$.
Przeliczalna rodzina takich kostek pokrywa $\Omega$,
więc suma wyjątkowych zbiorów miary zero również
ma miarę zero.
Jeśli $z\in H^2_{\mathrm{loc}}$, to dla każdego $j$
funkcja $\partial_jz\in H^1_{\mathrm{loc}}$.
Na mocy pierwszej części $E$ jest, z dokładnością
do zbioru miary zero, zawarty w zbiorze
$\{\partial_jz=0\}$. Ponownie stosując pierwszą część,
tym razem do $\partial_jz$, otrzymujemy
$\partial_k\partial_jz=0$ prawie wszędzie na $E$.
Biorąc skończoną sumę wyjątkowych zbiorów
dla wszystkich $j,k$, otrzymujemy tezę o $D^2z$.
\end{proof}

\begin{twierdzenie}[Regularność $H^2$ i nierówność Lewy'ego--Stampacchii]
\label{thm:pde29-LS}
Niech $\Omega\subset\R^d$ będzie ograniczonym obszarem
klasy $C^{1,1}$, $f\in L^2(\Omega)$ i
$\psi\in H^2(\Omega)$, przy czym ślad $\psi$ na
$\partial\Omega$ jest niedodatni. Niech $u\in H_0^1(\Omega)$
będzie rozwiązaniem zagadnienia przeszkody dla $A=I$.
Wówczas $u\in H^2(\Omega)$, a $\mu=-\Delta u-f$ jest funkcją
z $L^2(\Omega)$ i prawie wszędzie zachodzi
\begin{equation}\label{eq:pde29-LS}
 0\le\mu\le(-\Delta\psi-f)^+,
 \qquad f\le-\Delta u\le\max\{f,-\Delta\psi\}.
\end{equation}
Ponadto $(u-\psi)\mu=0$ prawie wszędzie w $\Omega$.
\end{twierdzenie}
\begin{proof}
\textbf{1. Zadanie z karą.}
Dla $\eps>0$ minimalizujemy na $V=H_0^1(\Omega)$
funkcjonał
\[
 J_\eps(v)=\frac12\|\nabla v\|_2^2
 -\int_\Omega fv\,\mathrm d\lambda_d(x)
 +\frac1{2\eps}\int_\Omega((\psi-v)^+)^2\,\mathrm d\lambda_d(x).
\]
Nieujemny składnik kary zachowuje koercywność
funkcjonału $J_f$. Funkcja $s\mapsto(s^+)^2$
jest wypukła (jej pochodna $2s^+$ jest niemalejąca),
więc całka kary jest wypukła. Energia Dirichleta
jest ściśle wypukła na $V$.
Dla ograniczonego ciągu minimalizującego wybieramy
podciąg słabo zbieżny w $V$ i silnie zbieżny w $L^2$
na mocy twierdzenia Rellicha~\ref{thm:pde10-Rellich-p}.
Ponieważ $s\mapsto s^+$ jest $1$-lipschitzowska,
$(\psi-v_n)^+\to(\psi-v)^+$ w $L^2$, więc kara
zbiega, a energia jest słabo półciągła z dołu.
Otrzymujemy jedyny minimizator $u_\eps$.

Dla $\varphi\in V$ różniczkujemy
$J_\eps(u_\eps+t\varphi)$ w $t=0$.
Funkcja $F(s)=(s^+)^2/2$ ma pochodną $F'(s)=s^+$
o stałej Lipschitza $1$. Z twierdzenia Newtona--Leibniza
wynika $|F(s+h)-F(s)-F'(s)h|\le |h|^2/2$.
Stosując to punktowo do $s=\psi-u_\eps$ i $h=-t\varphi$,
po całkowaniu otrzymujemy resztę ograniczoną przez
$t^2\|\varphi\|_2^2/(2\eps)$ dla składnika kary.
Zatem warunek stacjonarności daje dla każdego $\varphi\in V$
\[
 \int_\Omega\nabla u_\eps\cdot\nabla\varphi\,
 \mathrm d\lambda_d=\int_\Omega(f+\mu_\eps)\varphi\,
 \mathrm d\lambda_d,\qquad
 \mu_\eps=\eps^{-1}(\psi-u_\eps)^+\ge0.
\]
Jest to równanie $-\Delta u_\eps=f+\mu_\eps$ w $V'$.

\textbf{2. Jednostajna kontrola reakcji.}
Niech $w_\eps=(\psi-u_\eps)^+$.
Operator śladu komutuje z częścią dodatnią:
dla gładkich przybliżeń w $H^1$ ich części dodatnie
są ograniczone w $H^1$ i zbiegają w $L^2$,
a więc zbiegają słabo w $H^1$.
Z ciągłości śladu i jego zgodności z ograniczeniem
funkcji ciągłych (Twierdzenie~\ref{thm:pde10-trace})
otrzymujemy słabą zbieżność śladów w $L^2$
na brzegu; jednocześnie części dodatnie śladów
zbiegają silnie w $L^2$ dzięki własności Lipschitza.
Stąd ślad $w_\eps$ jest równy części dodatniej
śladu $\psi-u_\eps$, czyli zeru.
Na mocy Twierdzenia~\ref{thm:pde10-trace-zero}
mamy $w_\eps\in V$. Ten sam argument dla $\psi^+$
uzasadnia niepustość $K_\psi$ w hipotezach twierdzenia.
Wprowadzamy $h=-\Delta\psi-f\in L^2(\Omega)$.
Odejmujemy równanie $-\Delta\psi=f+h$ od
równania dla $u_\eps$ i testujemy różnicę przez
$w_\eps$. Reguła obcięcia
~\ref{cor:pde9-truncation} daje
$\nabla w_\eps=\chi_{\{\psi>u_\eps\}}
\nabla(\psi-u_\eps)$, więc lewa strona równania
dla różnicy wynosi $-\|\nabla w_\eps\|_2^2$.
Prawa strona jest równa
$\int_\Omega(\mu_\eps-h)w_\eps\,\mathrm d\lambda_d$.
Podstawiając $w_\eps=\eps\mu_\eps$ i przestawiając
składniki, dostajemy
\begin{equation}\label{eq:pde29-penalty}
 \|\nabla w_\eps\|_2^2+\eps\|\mu_\eps\|_2^2
 =\eps\int_\Omega h\mu_\eps\,\mathrm d\lambda_d(x)
 \le\eps\|h^+\|_2\|\mu_\eps\|_2.
\end{equation}
W szczególności $\|\mu_\eps\|_2\le\|h^+\|_2$,
$\|w_\eps\|_2\le\eps\|h^+\|_2$ oraz
$\|\nabla w_\eps\|_2^2\le
\eps\|h^+\|_2^2\to0$.
Testujemy równanie $-\Delta u_\eps=f+\mu_\eps$ przez $u_\eps$.
Z nierówności Poincarégo i oszacowania reakcji otrzymujemy
\[
 \|\nabla u_\eps\|_2^2
 =\int_\Omega(f+\mu_\eps)u_\eps\,\mathrm d\lambda_d(x)
 \le C_P(\|f\|_2+\|h^+\|_2)\|\nabla u_\eps\|_2.
\]
Zatem $\|u_\eps\|_V\le C_P(\|f\|_2+\|h^+\|_2)$
niezależnie od $\eps$, co uzasadnia wybranie słabo zbieżnego
podciągu w następnym kroku.

\textbf{3. Granica i warunek komplementarności.}
Dla dowolnego ciągu $\eps_n\downarrow0$
z jednostajnych oszacowań, refleksywności i twierdzenia
Rellicha wydzielamy podciąg (zachowujemy zapis $\eps$)
taki, że $u_\eps\rightharpoonup u$ w $V$,
$u_\eps\to u$ w $L^2(\Omega)$ i
$\mu_\eps\rightharpoonup\mu$ w $L^2(\Omega)$.
Nierówność Lipschitza daje
$(\psi-u_\eps)^+\to(\psi-u)^+$ w $L^2$,
a norma pierwszej funkcji dąży do zera, więc $u\ge\psi$.
Dla ustalonego $\varphi\in V$ przechodzimy
do granicy w równaniu: gradient zbiega słabo w $L^2$,
a reakcja słabo w $L^2$, gdzie $\varphi\in L^2$.
Otrzymujemy $-\Delta u=f+\mu$ w $V'$.
Dodatniość $\mu$ wynika z testowania słabej
zbieżności przez $\mu^-\in L^2$: mamy
$0\le\int_\Omega\mu\mu^-\,\mathrm d\lambda_d
=-\|\mu^-\|_2^2$.
Dodatkowo
\[
 \int_\Omega(u_\eps-\psi)\mu_\eps\,\mathrm d\lambda_d(x)
 =-\eps\|\mu_\eps\|_2^2\longrightarrow0.
\]
Różnicę całek rozbijamy na
$\int_\Omega(u_\eps-u)\mu_\eps\,\mathrm d\lambda_d$
oraz $\int_\Omega(u-\psi)(\mu_\eps-\mu)\,\mathrm d\lambda_d$.
Pierwsza dąży do zera z nierówności
Cauchy'ego--Schwarza i jednostajnego ograniczenia
$\|\mu_\eps\|_2$, a druga ze słabej zbieżności reakcji. Stąd
$\int_\Omega(u-\psi)\mu\,\mathrm d\lambda_d(x)=0$;
nieujemność czynników daje $(u-\psi)\mu=0$ prawie wszędzie.
Dla $v\in K_\psi$ zachodzi więc
$\int_\Omega\mu(v-u)\,\mathrm d\lambda_d(x)
 =\int_\Omega\mu(v-\psi)\,\mathrm d\lambda_d(x)\ge0$.
Wynika stąd~\eqref{eq:pde29-VI}, a jednoznaczność z
twierdzenia~\ref{thm:pde29-existence} identyfikuje $u$.
W istocie zachodzi silna zbieżność całej rodziny
w $V$, z oszacowaniem
\[
 \|u_\eps-u\|_V^2\le\tfrac{\eps}{4}\|\mu\|_2^2.
\]
Odejmujemy równania dla $u_\eps$ i $u$ oraz testujemy
przez $u_\eps-u$. Ponieważ $\mu_\eps\ge0$, $u\ge\psi$,
mamy $\mu_\eps(u_\eps-u)\le-\eps\mu_\eps^2$.
Z kolei komplementarność daje
$-\mu(u_\eps-u)=\mu(\psi-u_\eps)\le\eps\mu\mu_\eps$.
Po całkowaniu
$\|u_\eps-u\|_V^2\le
\eps(-\|\mu_\eps\|_2^2+(\mu,\mu_\eps)_{L^2})
\le\eps\|\mu\|_2^2/4$,
przez dopełnienie kwadratu w $L^2$.
Argument dotyczy każdego $\eps>0$, nie tylko podciągu.

\textbf{4. Regularność i nierówność punktowa.}
Z równania $-\Delta u=f+\mu\in L^2(\Omega)$ oraz
globalnego oszacowania regularności Dirichleta na obszarze
klasy $C^{1,1}$ (Twierdzenie~\ref{thm:pde15-boundary}
w rozdziale~\ref{chap:pde15}) otrzymujemy
$u\in H^2(\Omega)$ oraz
$\|u\|_{H^2}\le C_\Omega(\|f\|_2+\|h^+\|_2)$,
gdyż $\|\mu\|_2\le\|h^+\|_2$ ze słabej
półciągłości normy. Niech $z=u-\psi\in H^2(\Omega)$.
Z lematu~\ref{lem:pde29-zero-level} wynika
$D^2z=0$ prawie wszędzie na zbiorze $\{z=0\}$.
Na zbiorze kontaktu $\{u=\psi\}$ mamy więc
$-\Delta u=-\Delta\psi$ prawie wszędzie;
na jego dopełnieniu $\mu=0$ na mocy komplementarności.
Zatem $\mu=h$ prawie wszędzie na $\{u=\psi\}$,
gdzie dodatniość $\mu$ wymusza $h\ge0$,
i $\mu=0$ prawie wszędzie na $\{u>\psi\}$.
Otrzymujemy $0\le\mu\le h^+$ i po dodaniu $f$
drugi zapis~\eqref{eq:pde29-LS}.
\end{proof}

\begin{uwaga}[Granica regularności]
Oszacowanie $H^2$ zależy od regularności brzegu, a
nierówność~\eqref{eq:pde29-LS} od drugich pochodnych przeszkody.
Nie wolno ich przypisywać rozwiązaniu dla dowolnej
przeszkody $H^1$ ani dowolnego obszaru Lipschitzowskiego.
Na ogólnych obszarach $\mu$ jest na początku jedynie miarą
Radona, a swobodny brzeg może być nieregularny.
\end{uwaga}

\section{Zbiór kontaktu i swobodny brzeg}
Przyjmijmy założenia Twierdzenia~\ref{thm:pde29-LS}
oraz $d\le3$. Wówczas $u,\psi\in H^2(\Omega)$
mają jednoznaczne ciągłe reprezentanty na $\overline\Omega$.
Istotnie, dla $d=3$ stosujemy Twierdzenie
~\ref{thm:pde10-domain-embed} do funkcji i jej pochodnych
pierwszego rzędu: $H^2\hookrightarrow W^{1,6}$,
a następnie $W^{1,6}\hookrightarrow C^{0,1/2}$.
Dla $d=2$ analogicznie wybieramy wykładnik $4$,
otrzymując $H^2\hookrightarrow W^{1,4}\hookrightarrow C^{0,1/2}$.
Dla $d=1$ wystarcza jednowymiarowa charakterystyka
$H^1$ z Twierdzenia~\ref{thm:pde9-ac}.
Ciągły reprezentant różnicy $u-\psi$ jest nieujemny
w każdym punkcie $\Omega$: gdyby był ujemny w pewnym
punkcie, byłby ujemny w całej kuli dodatniej miary,
wbrew warunkowi dopuszczalności.
Definiujemy wtedy
$\mathcal C=\{x\in\Omega:u(x)=\psi(x)\}$ (zbiór kontaktu)
oraz $\mathcal N=\{x\in\Omega:u(x)>\psi(x)\}$ (zbiór braku kontaktu).
Zbiory $\mathcal C$ i $\mathcal N$ są odpowiednio
względnie domknięty i otwarty w $\Omega$, a
$\Omega=\mathcal C\mathbin{\dot\cup}\mathcal N$.
Zbiorem swobodnego brzegu nazywamy
$\Gamma=\partial_\Omega\mathcal N
=\mathcal C\cap\overline{\mathcal N}^{\,\Omega}$,
gdzie domknięcie i brzeg rozumiemy w topologii względnej.
Z komplementarności wynika $\mu=0$ prawie wszędzie
na $\mathcal N$, więc $-\Delta u=f$ w sensie
dystrybucji na tym zbiorze. Można to też wykazać
bez korzystania z regularności reakcji: dla
$\varphi\in C_c^\infty(\mathcal N)$ różnica
$u-\psi$ ma dodatnie minimum na nośniku $\varphi$
(o ile $\varphi\ne0$), zatem oba testy
$u\pm t\varphi$ są dopuszczalne dla dostatecznie małego
$t>0$. Nierówność wariacyjna dla obu znaków
daje $\int_\Omega\nabla u\cdot\nabla\varphi\,
\mathrm d\lambda_d=\int_\Omega f\varphi\,\mathrm d\lambda_d$.
Na $\mathcal C$ z Lematu~\ref{lem:pde29-zero-level}
wynika $-\Delta u=-\Delta\psi$ prawie wszędzie,
a zatem $\mu=(-\Delta\psi-f)\chi_{\mathcal C}$
prawie wszędzie. Nie utożsamiamy zbioru kontaktu
z nośnikiem reakcji: reakcja może znikać na części
zbioru kontaktu (na przykład dla $\psi=f=0$). Zbiór $\Gamma$ nie jest dany
z góry i jego regularność wymaga dodatkowej analizy.

\begin{przyklad}[Swobodny brzeg i skok drugiej pochodnej w wymiarze jeden]
\label{ex:pde29-explicit}
Niech $\Omega=(-1,1)$, $f=0$ i
$\psi(x)=(1-x^2)^2$.
Szukamy parzystego rozwiązania, które na $[0,a]$
pokrywa się z przeszkodą, a na $[a,1]$ jest afiniczne
i znika w punkcie $1$, gdzie $0<a<1$.
Ciągłość wymusza wzór $u(x)=m(1-x)$ na $[a,1]$,
z $m=\psi(a)/(1-a)=(1-a)(1+a)^2$.
Ciągłość pochodnej w $a$ wymaga
$m=-\psi'(a)=4a(1-a^2)$.
Po podzieleniu przez $(1-a)(1+a)>0$
otrzymujemy $1+a=4a$, czyli $a=1/3$ i $m=32/27$.
Jest to sposób znalezienia kandydata, a nie
założenie o postaci rozwiązania.
Zweryfikujemy, że jedynym rozwiązaniem jest funkcja
\begin{equation}\label{eq:pde29-explicit}
 u(x)=\begin{cases}
 (1-x^2)^2,& |x|\le\frac13,\\
 \frac{32}{27}(1-|x|),&\frac13<|x|\le1.
 \end{cases}
\end{equation}
Istotnie, $u(\pm1)=0$. W punkcie $1/3$
obie wartości wynoszą $64/81$, a obie pochodne
$-32/27$; w punkcie $-1/3$ korzystamy z parzystości.
Funkcja jest więc klasy $C^1$ na $[-1,1]$,
a jej pochodna jest odcinkami gładka z ograniczoną
drugą pochodną. Stąd $u'\in W^{1,\infty}(-1,1)$,
gdyż całkowanie przez części na trzech podprzedziałach
nie daje wyrazów skokowych w punktach sklejenia.
W szczególności $u\in H^2(-1,1)\cap H_0^1(-1,1)$. Dla $x\in[1/3,1]$ zachodzi
\[
 u(x)-\psi(x)=\frac{(1-x)(3x-1)^2(3x+5)}{27}\ge0;
\]
na lewej połowie korzystamy z parzystości.
Dla $1/3<x<1$ wszystkie czynniki są dodatnie,
więc zbiór kontaktu wewnątrz $\Omega$ jest dokładnie
$[-1/3,1/3]$.
Ponadto $-u''=(4-12x^2)\chi_{\{|x|<1/3\}}\ge0$
prawie wszędzie. Jest to również dystrybucyjna
druga pochodna ze zmienionym znakiem,
bez mas Diraca w punktach sklejenia.
Reakcja jest zerowa poza zbiorem kontaktu,
a na nim $v-u=v-\psi\ge0$ dla $v\in K_\psi$.
Całkowanie przez części (z gęstości funkcji
testowych w $H_0^1$) daje więc
$\int_{-1}^1u'(v'-u')\,\mathrm d\lambda_1(x)
=\int_{-1}^1(-u'')(v-u)\,\mathrm d\lambda_1(x)\ge0$.
Z jednoznaczności wynika wzór~\eqref{eq:pde29-explicit}.
Swobodny brzeg składa się z dwóch punktów
$\Gamma=\{-1/3,1/3\}$. Funkcja $u$ należy do
$C^{1,1}([-1,1])$, ale nie do $C^2([-1,1])$,
ponieważ jednostronne granice $u''$ w punkcie $1/3$
wynoszą $-8/3$ od lewej i $0$ od prawej.
W punkcie $-1/3$ kolejność jest odwrotna.
Całkowita masa reakcji wynosi
\[
 \int_{-1}^1\mu\,\mathrm d\lambda_1
 =2\int_0^{1/3}(4-12x^2)\,\mathrm d\lambda_1(x)
 =\tfrac{64}{27}.
\]
Jest ona równa $u'(-1)-u'(1)$, zgodnie
z całkowaniem równania $-u''=\mu$.
\end{przyklad}

\begin{center}
\begin{tikzpicture}[x=4.4cm,y=2.5cm]
 \draw[-{Latex}] (-1.06,0)--(1.08,0) node[right] {$x$};
 \draw[-{Latex}] (-1,-.12)--(-1,1.15) node[above] {$u,\psi$};
 \draw[dashed,gray] (-.3333,0)--(-.3333,1.07);
 \draw[dashed,gray] (.3333,0)--(.3333,1.07);
 \fill[blue!6] ({-1/3},0) rectangle ({1/3},1);
 \draw[blue!70!black,thick,samples=181,domain=-1:1]
   plot (\x,{(1-\x*\x)^2});
 \draw[red!75!black,very thick,samples=81,domain=-1/3:1/3]
   plot (\x,{(1-\x*\x)^2});
 \draw[red!75!black,very thick] (-1,0)--({-1/3},{64/81})
   ({1/3},{64/81})--(1,0);
 \node[below] at (-1,0) {$-1$};
 \node[below] at (-.3333,0) {$-\tfrac13$};
 \node[below] at (.3333,0) {$\tfrac13$};
 \node[below] at (1,0) {$1$};
 \node[blue!65!black,right] at (.58,.26) {$\psi$};
 \node[red!70!black,right] at (.62,.54) {$u$};
 \node[above] at (0,1.02) {zbiór kontaktu};
\end{tikzpicture}

{\small Rysunek: przeszkoda (niebieska) i rozwiązanie (czerwone)
z Przykładu~\ref{ex:pde29-explicit}. W środku zachodzi $u=\psi$;
poza tym przedziałem $-u''=0$.}
\end{center}

\begin{center}
\begin{tikzpicture}[x=4.4cm,y=.66cm,>=Latex,font=\small]
 \fill[orange!20] ({-1/3},0)
  -- plot[domain=-1/3:1/3,samples=81] (\x,{4-12*\x*\x})
  -- ({1/3},0)--cycle;
 \draw[->] (-1.06,0)--(1.08,0) node[right] {$x$};
 \draw[->] (-1,-.12)--(-1,4.7) node[above] {$\mu$};
 \draw[orange!85!black,very thick] (-1,0)--({-1/3},0)
    ({1/3},0)--(1,0);
 \draw[orange!85!black,very thick]
    plot[domain=-1/3:1/3,samples=81] (\x,{4-12*\x*\x});
 \foreach \a in {-1/3,1/3}{
   \draw[gray,dashed] ({\a},0)--({\a},{8/3});
   \draw[orange!85!black,fill=white] ({\a},0) circle[radius=1.5pt];
   \draw[orange!85!black,fill=white] ({\a},{8/3}) circle[radius=1.5pt];
 }
 \foreach \a/\lab in {-1/-1,-1/3/{-\frac13},1/3/{\frac13},1/1}
   \node[below] at ({\a},0) {$\lab$};
 \draw (-1,4)--++(-.025,0) node[left] {$4$};
 \draw (-1,{8/3})--++(-.025,0) node[left] {$8/3$};
 \node[above] at (0,4.15) {reakcja na zbiorze kontaktu};
\end{tikzpicture}

{\small Rysunek: gęstość reakcji $\mu=-u''$.
Przerywane odcinki zaznaczają skoki; wartości w pojedynczych
punktach swobodnego brzegu nie mają znaczenia dla klasy $L^2$.}
\end{center}

\section{Projekcja na zbiór wypukły i mnożniki ograniczeń}
Dla funkcjonału $f\in V'$ niech $w\in V$ będzie jedynym
rozwiązaniem $a(w,v)=\langle f,v\rangle$ dla $v\in V$
(twierdzenie Laxa--Milgrama). Wówczas
$J_f(v)=\frac12a(v-w,v-w)-\frac12a(w,w)$.
Rozwiązanie nierówności wariacyjnej jest zatem rzutem $w$
na $K_\psi$ w metryce określonej przez iloczyn skalarny $a$.
Nie jest to na ogół rzut punktowy $\max\{w,\psi\}$:
energia kontroluje pochodne, więc zmiana w jednym miejscu
wpływa na rozwiązanie w całym obszarze.

Dla domkniętego wypukłego zbioru $K\subset V$ definiujemy
stożek normalny
$N_K(u)=\{\eta\in V':\langle\eta,v-u\rangle\le0
\text{ dla każdego }v\in K\}$, jeżeli $u\in K$,
i $N_K(u)=\varnothing$ w przeciwnym przypadku.
Jeżeli $\mathcal A:V\to V'$ jest operatorem odpowiadającym $a$,
to nierówność~\eqref{eq:pde29-VI} jest równoważna
$0\in\mathcal Au-f+N_{K_\psi}(u)$.
Ta postać pozwala przenieść argumenty na inne problemy
jednostronne, także na przeszkody ruchome oraz przepływy
gradientowe z ograniczeniami.

\section{Schemat numeryczny i warunki punktowe}
W wymiarze jeden, dla $\Omega=(-1,1)$ i równomiernej siatki
$x_i=-1+ih$, $h=2/N$, $i=0,\ldots,N$, oznaczamy
$U_i\approx u(x_i)$ i $\Psi_i=\psi(x_i)$.
Dla $f=0$ dyskretna energia jest równa
$E_h(U)=\frac1{2h}\sum_{i=0}^{N-1}(U_{i+1}-U_i)^2$
przy $U_0=U_N=0$. Minimalizujemy ją na zbiorze
$U_i\ge\Psi_i$ dla $i=1,\dots,N-1$.
W węzłach wewnętrznych warunki optymalności mają postać
\begin{equation}\label{eq:pde29-discrete}
 U_i\ge\Psi_i,\qquad
 (2U_i-U_{i-1}-U_{i+1})/h^2\ge0,\qquad
 (U_i-\Psi_i)(2U_i-U_{i-1}-U_{i+1})=0.
\end{equation}
Dla przeszkody z przykładu~\ref{ex:pde29-explicit}
można zastosować metodę Gaussa--Seidela z projekcją,
zastępując każdy nowy węzeł przez większą z wartości
$\Psi_i$ i średniej sąsiednich wartości. Sprawdzanie
nierówności~\eqref{eq:pde29-discrete} jest ważniejsze niż
sam obraz krzywej: wykres może wyglądać poprawnie również
wtedy, gdy warunek komplementarności jest naruszony.

\section{Dalsze pytania i granice wykazanej teorii}
Udowodniliśmy istnienie i jednoznaczność nierówności
wariacyjnej w $H_0^1$, a przy mocniejszych założeniach
regularność $H^2$ i nierówność Lewy'ego--Stampacchii.
Nie udowodniliśmy ogólnej regularności $C^{1,1}$
przy dowolnej gładkiej przeszkodzie ani regularności
swobodnego brzegu. Analiza jego punktów regularnych
i osobliwych wymaga dodatkowych oszacowań wzrostu,
monotoniczności i argumentów geometrycznych;
nie wynika z samego sformułowania wariacyjnego.


